\section{AW1: parity, link flip and the first-order Wilson mean}
\label{sec:aw1}

\textbf{Contribution \projecttag{HNM-C-AW1}.} The AW1 contract, recorded under
the title ``Hruday parity theorem, link-flip antisymmetry and the first-order
Wilson mean with an itemized second-order remainder'', was frozen after the
AV2 gate and before production. It keeps the model of
Sections~\ref{sec:av1} and~\ref{sec:av2}: the zero selected triple, whole
stars under I1.5, the cover $R=\{0,e_z\}$ and both signs of
$|\tau|\le10^{-8}$. Its preregistered target is $K_2^+\tau\le1/288$ at the
cap, half the first-order coefficient. It also freezes the coupling rule for
AW2 before any constant is known. Both directions are paired. Beyond the
shared premises, the forward producer received the skeptic's triage, the
skeptic's loop-2 response and the historical and modern loop-2 responses.
These already state the parity rule, the flip set and the rough structure of
$K_2$, so they are disclosed as shared panel premises for that route. The
reverse producer received only \texttt{AGENTS.md}, the contract and its 29
shared premises, 31 files in all. It reconstructs the parity rule and the flip
set from the character structure.

Notation: $\mathbb E=E_{\rm Haar}$ is the Haar expectation of the reference,
$E$ is the flip set, and $E_N$ is the ground energy of the box $\Lambda_N$. In
units of $\alpha$ the finite-box Hamiltonian is $G_N=H_N/\alpha$ and the
omitted interaction is $V=\sum_bV_b$ with $V_b=-(\tau/24)\sum_{f\in O_b}W_f$;
$\kappa=(\lambda_L,\mu_M,\lambda_R)$ is the selected triple. Boxes are open
centered whole-star boxes, with no periodic identification.

\subsection{Re-targeting roadmap goal 2}
\label{sec:aw1-retarget}
The second Round31 roadmap goal asks for ``a certified quantity whose
interacting enclosure excludes its matched free reference''
(\repo{research/round31/advisor/roadmap.json}{Round31 roadmap}). AV2 did not
meet it: the certified interval for $C_\tau(1)$ contains $e^{-3}/4$
(Section~\ref{sec:av2-certificate}). Before AW1 was contracted, the panel
observed in deliberation that SU(2) Peter--Weyl parity removes every
first-order term of the centered correlation in the zero-selected subfamily,
so that its interaction shift is of second order in $\tau$. AW1 proves this
for finite boxes, with a constant not uniform in the volume
(Section~\ref{sec:aw1-parity}). A certificate whose error is linear in
$\tau$, such as AV2's, cannot resolve a shift of second order.

The contract therefore re-targets goal~2 to the only first-order Wilson effect
in this model, the uncentered equal-time mean $\omega_\tau(W)$. Its candidate
expansion $+\tau/144+O(\tau^2)$ was to be derived, not assumed. A sign
certificate for the mean needs three ingredients: the exact coefficient, an
exact antisymmetry control and a second-order remainder uniform in volume,
taken from the same AM2 expansion admitted through AV1. AW1 supplies the three
ingredients and fixes the coupling of the certificate. The certificate itself
belongs to AW2 (Section~\ref{sec:aw2}). The mean is a static ground-state
quantity. Neither its value nor its sign is a dynamical, spectral or mass-gap
correction, and the shift of the centered correlation remains unresolved.

\subsection{Reviewed scope}
The \repo{research/round32/advisor/aw1-gate.json}{AW1 gate} records the
verdict \recid{accepted_within_scope}. The skeptic's review carries the
sub-label \recid{static_not_dynamic}, and the gate's statement and limitations
carry the same restriction. The accepted statement, quoted verbatim below in
the gate's ASCII notation, is the scope of every claim in this section. In it,
\texttt{E[$\cdot$]} is a Haar expectation, the \texttt{E} in \texttt{G-E} is
the ground energy, and \texttt{E} alone is the flip set.
\begin{gatequote}
In the AM2/AQ1 zero-selected patterned family (SU(2) Kogut-Susskind form on
Z\textasciicircum{}3 at fixed spacing, coarse 24-link factors, selected triple
exactly (0,0,0) with Haar product reference, whole stars phi\_b=-(tau/3) sum
W\_f in normalized units delta=alpha/8; original xz Wilson loop W with
complete cover R=\{0,e\_z\}; clock s=alpha*t\_E/hbar, G=H/alpha), at both
signs of tau with |tau|<=10\textasciicircum{}-8: (1) Parity theorem, finite
boxes only: in every open centered whole-star box Lambda\_N, N>=2, and every
on-site cutoff, the first-order tau-derivatives at tau=0 of
omega\_N(W\textasciicircum{}2), of the real-time correlation
c\_N(theta)=<chi,e\textasciicircum{}\{i theta(G-E)\}chi> (every real theta)
and of C\_N(s)=<chi,e\textasciicircum{}\{-s(G-E)\}chi> (every s>=0) vanish,
the state, Duhamel, energy and vector-centring terms separately, by the
per-link Z\_2 centre grading (SU(2) Peter-Weyl multiplicities):
E[W]=E[W\textasciicircum{}3]=0, E[W\textasciicircum{}2]=1/4,
E[W\textasciicircum{}4]=1/8, E[W\textasciicircum{}2 W\_f]=0 for every omitted
f including f=W, P\_24 V W Omega\_0=0 (only the identity component
(1/4)Omega\_0 at energy 0 for f=W), and zero first-order splitting of the
gauge-invariant energy-24 multiplet span\{W\_g Omega\_0\} (not claimed on the
non-invariant eigenspace); hence
C\_N(s)=e\textasciicircum{}\{-3s\}/4+O\_N(tau\textasciicircum{}2) with the
tau\textasciicircum{}2 constant explicitly unbounded (not uniform in N; for AQ
limits only evenness of C(s), no AQ second-order statement), and
omega(W\textasciicircum{}2)=1/4+O(tau\textasciicircum{}2) with supplementary
uniform constants (not targets; the free value lies inside, no shift
resolved). (2) Flip lemma: E=\{(p,x):p\_y even\} u \{(p,y):p\_z even\} u
\{(p,z):p\_x even\} meets every plaquette of Z\textasciicircum{}3 in 1 or 3
links; U\_E (central -1 on every link of E) commutes with every Casimir and
every original endpoint gauge action, fixes Omega\_0 and maps every W\_f to
-W\_f, so U\_E H\_N(tau,kappa) U\_E\textasciicircum{}*=H\_N(-tau,-kappa) in
every open centered whole-star box, including on-site cutoff compressions
(Q\_L -> cutoff of h\_b(-kappa)); at kappa=0
omega\_\{N,-tau\}(W)=-omega\_\{N,tau\}(W) and omega\_N(W\textasciicircum{}2),
C\_N, c\_N are even in tau; for AQ limits S(-tau)=S(tau)o alpha\_E as whole
sets of subsequential limits (pointwise only along a common subsequence); odd
periodic sides and frozen or gauge-fixed links excluded; no tau-antisymmetry
via U\_E at kappa!=0; oddness gives no O(tau\textasciicircum{}3) remainder.
(3) omega\_tau(W)=+tau/144+r(tau) under I1.5 (AV1 creation convention
c\textasciicircum{}(1)=L\_0=-(tau/72) sum W\_f Omega\_0, omega=-2Re<W
Omega\_0,c\textasciicircum{}(1)>+O(tau\textasciicircum{}2)); only f=W
contributes; first-order values +-1/14400000000 at
tau=+-10\textasciicircum{}-8. (4) |r(tau)|<=K\_2\textasciicircum{}+
tau\textasciicircum{}2 uniformly in N, cutoff and every AQ1 subsequential
limit at the exact tier (t<=T=(49|tau|/144)/(1-352J), rho=352JT, J=28|tau|):
the bound value is
K\_2\textasciicircum{}+=81108864767825329926713064490531229475390625/\allowbreak
24176936535511801466930759024724079017984
(\textasciitilde{}3354.80322946, outward
3354803229461691/10\textasciicircum{}12; the skeptic itemization, identical to
the forward\textquotesingle{}s labelled unpinned\_t\_bounds variant,
dominating the forward 3354.63832966 and reverse 3354.49942587 exact-tier
values, all three valid). (5) K\_2\textasciicircum{}+ tau<=1/288 at
tau=10\textasciicircum{}-8, so the frozen decade-grid rule gives
tau\_AW2=10\textasciicircum{}-8 (the cap) with sign margin 1/(144
K\_2\textasciicircum{}+ tau)\textasciitilde{}207.00005 (directed floor
41400010489/200000000). The crude tier (t\_c=592|tau|,
K\_2\textasciitilde{}7.9375e6, margin 0.0875) fails at the cap and is retained
as a limited-tier value. No enclosure or sign of omega(W), no dynamical or
mass-gap statement, no third-order remainder, uniqueness, whole-sequence
convergence, rate in N, uniform Wilson, weak-coupling, continuum or priority
claim is admitted.
\end{gatequote}
The rest of this section derives that statement. The limitations recorded by
the gate close the section (Section~\ref{sec:aw1-limits}).

\subsection{The parity theorem}
\label{sec:aw1-parity}
\paragraph{Centre grading.}
For a link $e$ let $(\Gamma_e\psi)(\ldots,U_e,\ldots)=\psi(\ldots,-U_e,\ldots)$.
Haar measure is invariant under the central element $-1$, so $\Gamma_e$ is
unitary and fixes $\Omega_0$. On the Peter--Weyl block $V_j\otimes V_j^*$ of
the link it acts as $(-1)^{2j}$, while the Casimir $C_e$ acts as $j(j+1)$. The
grading therefore commutes with every Casimir, with $H_0$ and with every
function of the Casimirs: $H_0^{-1}P_\perp$, $e^{-sH_0}$, $e^{i\theta H_0}$
and every on-site spectral cutoff at the zero triple. Because $-1$ is central,
$\Gamma_e$ also commutes with left and right translations, hence with every
original endpoint gauge action. Finally,
$\Gamma_eW_f\Gamma_e=(-1)^{[e\in f]}W_f$. In character form, the Haar integral
over one link of $k$ spin-$\tfrac12$ matrix elements projects onto the
invariants of $V_{1/2}^{\otimes k}$. Their multiplicity is zero for odd $k$
and the Catalan number of $k/2$ for even $k$. Hence, for centre-even $A_i$
such as functions of the Casimirs,
\begin{equation}
 \langle\Omega_0,A_0W_{f_1}A_1\cdots W_{f_k}A_k\Omega_0\rangle=0
 \quad\text{unless every link lies in an even number of }f_1,\dots,f_k .
 \label{eq:aw1-parity}
\end{equation}

\paragraph{Haar moments and pairings.}
For one plaquette, $\mathbb E[W^n]$ is $2^{-n}$ times the multiplicity of spin
$0$ in $V_{1/2}^{\otimes n}$:
\begin{equation}
 \mathbb E[W^n]_{n=0,\dots,8}=1,\,0,\,\tfrac14,\,0,\,\tfrac18,\,0,\,\tfrac5{64},\,0,\,\tfrac7{128}.
 \label{eq:aw1-moments}
\end{equation}
Three exact routes agree through $n=8$: Clebsch--Gordan recursion, Weyl
integration by Wallis ratios, and the contract's named cross-check. For the
cross-check, the Wilson face contains the free $z$ link $(e_x,z)$. Under the
reference its holonomy is therefore Haar, $W=q_0$ with $q$ uniform on $S^3$,
and $\mathbb E[q_0^{2m}]=\prod_{i<m}(2i+1)/(4+2i)$. This confirms
$\mathbb E[W]=0$ and $\mathbb E[W^2]=\tfrac14$ without the character count.
Distinct plaquettes share at most one link. So \eqref{eq:aw1-parity} gives
$\mathbb E[W_f]=0$ and $\mathbb E[WW_f]=\tfrac14\delta_{f,W}$. It also gives
$\mathbb E[W^2W_f]=0$ for every omitted $f$, \emph{including} $f=W$: for
$f\ne W$ the links of $f$ outside $W$ carry one spin-$\tfrac12$ factor each,
and for $f=W$ every link of $W$ carries three. Finally
$\mathbb E[W_gW_fW_h]=0$ for all plaquettes $g,f,h$. The reverse proof uses
the flip set of Section~\ref{sec:aw1-flip}: $E$ meets each plaquette an odd
number of times, so it meets the $\mathbb Z_2$ chain $g+f+h$ an odd number of
times, and that chain is not empty. The forward proof notes that $g\triangle h$
is never a single face.

\paragraph{First-order terms.}
In each box the interaction is bounded and the ground energy is simple and
isolated (AM2). Kato's analytic perturbation theory (cited, not
machine-checked) then makes $\omega_N(W^2)$,
$c_N(\theta)=\langle\chi,e^{i\theta(G_N-E_N)}\chi\rangle$ and
$C_N(s)=\langle\chi,e^{-s(G_N-E_N)}\chi\rangle$ real-analytic in $\tau$ near
$0$. Write $\psi^{(1)}$ for the first-order correction of the ground vector.
AM2 writes $\psi=e^{-C}\Omega_0$, so with the first-order collection
\eqref{eq:av1-first},
\begin{equation}
 \psi^{(1)}=-c^{(1)}=+\frac{\tau}{72}\sum_fW_f\Omega_0 .
 \label{eq:aw1-psi1}
\end{equation}
The vector $W\Omega_0$ is an eigenvector of $G_0$ with eigenvalue $3$. The
first-order part of $C_N(s)$ is the sum of four terms, and each vanishes
separately:
\begin{equation}
 \begin{aligned}
 \text{state:}&\quad 2e^{-3s}\,\mathrm{Re}\langle W^2\Omega_0,\psi^{(1)}\rangle
 =\frac{\tau e^{-3s}}{36}\sum_f\mathbb E[W^2W_f]=0,\\
 \text{Duhamel:}&\quad -s\,e^{-3s}\langle W\Omega_0,(V-E^{(1)})\,W\Omega_0\rangle
 =\frac{\tau s e^{-3s}}{24}\sum_f\mathbb E[W^2W_f]=0,\\
 \text{energy:}&\quad E^{(1)}=\langle\Omega_0,V\Omega_0\rangle
 =-\frac{\tau}{24}\sum_f\mathbb E[W_f]=0 .
 \end{aligned}
 \label{eq:aw1-terms}
\end{equation}
The fourth, vector-centring, term needs care. The uncentered mean has the
nonzero first-order coefficient $1/144$ (Section~\ref{sec:aw1-mean}). It drops
out of the correlation only because it multiplies overlaps of $W\Omega_0$ with
$\Omega_0$, which equal $\mathbb E[W]=0$. Equivalently, since
$e^{-s(G_N-E_N)}\psi_N=\psi_N$, one has
$C_N(s)=\langle W\psi_N,e^{-s(G_N-E_N)}W\psi_N\rangle-m_N^2$ for the normalized
ground vector, and $m_N^2$ is of second order. The real-time correlation has
the same terms with $e^{3i\theta}$ and $i\theta e^{3i\theta}$ in place of
$e^{-3s}$ and $-se^{-3s}$. Its state term also pairs $c^{(1)}$ with
$W\alpha^0_\theta(W)\Omega_0=e^{3i\theta}W^2\Omega_0$, which is centre-even,
so the pairing vanishes. For the second moment the first-order part is
$2\mathrm{Re}\langle(W^2-\tfrac14)\Omega_0,\psi^{(1)}\rangle
=\tfrac{\tau}{36}\sum_f\bigl(\mathbb E[W^2W_f]-\tfrac14\mathbb E[W_f]\bigr)=0$.
Both checkers evaluate every sum exactly, over all 1344 retained faces of
$\Lambda_2$ and over the 82 faces meeting $R$.

\paragraph{The energy-24 multiplet.}
In normalized units the per-link content of $W_fW\Omega_0$ fixes its energy.
For $f=W$ every link carries $j\in\{0,1\}$, so the energy lies in
$\{0,16,32,48,64\}$. A face sharing one link with $W$ gives $\{36,52\}$, and a
disjoint face gives $48$. Hence
\begin{equation}
 P_{24}\,V\,W\Omega_0=0 .
 \label{eq:aw1-p24}
\end{equation}
The only component with the vacuum's content is the identity component
$\mathbb E[W^2]\Omega_0=\tfrac14\Omega_0$ of the $f=W$ term, at energy $0$. On
$\Lambda_2$ the odd-link sets $W\triangle f$ have sizes $0$, $6$ and $8$, for
$1$, $10$ and $1333$ omitted faces; the size $4$ of an energy-24 vector never
occurs. For the multiplet itself, $8\sum_ej_e(j_e+1)=24$ forces exactly four
spin-$\tfrac12$ links. Gauge invariance needs even degree at every vertex, so
the four links form a 4-cycle, and the 4-cycles of $\mathbb Z^3$ are exactly
the plaquettes. Each vertex of degree two carries a one-dimensional
invariant. The gauge-invariant energy-24 multiplet is therefore
$\mathrm{span}\{W_g\Omega_0\}$. On it
$\langle W_g\Omega_0,VW_h\Omega_0\rangle=-\tfrac{\tau}{24}\sum_f\mathbb E[W_gW_fW_h]=0$,
so its first-order splitting is zero. The statement is \emph{not} made for the
full energy-24 eigenspace. There, four spin-$\tfrac12$ links with two on a
face $f$ can be mapped back into the eigenspace by $W_f$. Both producers
restrict the claim to the gauge-invariant multiplet. That is where $\chi$ and
its dynamics live, because $G_N$ commutes with the gauge actions.

\paragraph{Exact scope of the conclusion.}
For every open centered whole-star box $\Lambda_N$, $N\ge2$, and every
on-site cutoff,
\begin{equation}
 C_N(s)=\tfrac14e^{-3s}+O_N(\tau^2),\qquad
 c_N(\theta)=\tfrac14e^{3i\theta}+O_N(\tau^2),\qquad
 \omega_N(W^2)=\tfrac14+O(\tau^2).
 \label{eq:aw1-parity-conclusion}
\end{equation}
Four qualifications are part of the admitted statement.
\begin{itemize}
\item \emph{Finite boxes.} The first two expansions are Taylor statements in
each box, from Kato analyticity per box.
\item \emph{Constant not uniform.} The $\tau^2$ constant of $C_N$ and $c_N$ is
explicitly unbounded. It is not proved uniform in $N$, so no second-order
statement about $C(s)$ or $c(\theta)$ passes to AQ limits. A uniform constant
would need a second-order dynamical estimate that neither route supplies.
\item \emph{Evenness in AQ limits.} For AQ limits only the evenness of $C(s)$
and $c(\theta)$ in $\tau$ passes, as a set-level statement from the flip
lemma (Section~\ref{sec:aw1-flip}).
\item \emph{Gauge-invariant multiplet.} Zero first-order splitting is claimed
for $\mathrm{span}\{W_g\Omega_0\}$ only.
\end{itemize}
For $\omega(W^2)$ the product split of Section~\ref{sec:av1-split} also gives
supplementary constants that are uniform in $N$, the cutoff and every AQ
subsequential limit: $|\omega(W^2)-\tfrac14|\le K_{W^2}\tau^2$ with $K_{W^2}$
about $3354.567$ (forward), $1677.331$ (reverse) and $1677.518$ (skeptic). The
forward value is valid but conservative, because it charges $2\varrho$ and all
72 straddling faces. These constants are not targets. The free value
$\tfrac14$ lies inside each band, so no shift is resolved.

\subsection{The link-flip lemma}
\label{sec:aw1-flip}
\paragraph{Odd intersection.}
The flip set, read from the contract, is
\begin{equation}
 E=\{(p,x):p_y\ \text{even}\}\cup\{(p,y):p_z\ \text{even}\}\cup\{(p,z):p_x\ \text{even}\}.
 \label{eq:aw1-flipset}
\end{equation}
Each link direction is keyed to the parity of the next coordinate in cyclic
order. In a face with directions $a<c$, the two $a$-links differ only in the
$c$ coordinate and the two $c$-links only in the $a$ coordinate. In each of
the three orientations exactly one of the two pairs is keyed to the
coordinate in which its links differ. That pair contributes exactly one link
of $E$, and the other pair contributes none or two. Hence $|f\cap E|\in\{1,3\}$
for every plaquette $f$ of $\mathbb Z^3$. Translation by an even vector
preserves $E$, so $|f\cap E|$ depends only on the orientation and on
$p\bmod2$: 24 classes. Both producers enumerate them in full, as the contract
requires. The forward producer also checks the 49 omitted faces of seven
factors at both $z$-parities, their selected faces and every plaquette of the
$N=2$ box. The reverse producer checks every plaquette with base point in
$[-9,9]^3$: 20\,577 plaquettes, 10\,830 meeting $E$ once, 9\,747 three times
and none evenly. It also checks all 49 omitted faces of the factors $0$ and
$e_z$. These two factors exhaust the $z$-parity classes, because coarse $x$
and $y$ translations are even fine shifts (4 and 2) while a coarse $z$
translation is a unit shift. The Wilson face meets $E$ in three links.

The reverse producer reconstructs $E$ before parsing it. Consider the 27
rules ``$(p,d)\in E$ iff $p_{\sigma(d)}$ is even'' for maps
$\sigma:\{x,y,z\}\to\{x,y,z\}$. Exactly two of them meet every plaquette
oddly: the cyclic set \eqref{eq:aw1-flipset} and the anti-cyclic set
$E'=\{(p,x):p_z\ \text{even}\}\cup\{(p,y):p_x\ \text{even}\}\cup\{(p,z):p_y\ \text{even}\}$.
They differ by the coboundary of $p_xp_y+p_yp_z+p_zp_x\bmod2$, which is a
centre gauge transformation. A solution can exist only if the constant
$\mathbb Z_2$ 2-cochain $1$ is a cocycle. It is, because a cube has six faces.
The set $E$ is not itself a cocycle, since its coboundary is $1$ on every
plaquette. So $U_E$ is not a gauge transformation, and it genuinely flips
gauge-invariant observables; the solution is unique up to centre gauge
transformations. Four mutations are rejected: $E$ minus one link, which leaves
four plaquettes even; an odd periodic side, which leaves even plaquettes at
the seam; a centre gauge transformation presented as the flip, which is even
on every plaquette; and a spin-one face term claimed to flip, since
$(-1)^{2j|f\cap E|}=+1$ for $j=1$.

\paragraph{The operator identity.}
Let $U_E=\prod_{e\in E\cap\Lambda_N}\Gamma_e$. It is diagonal in the
Peter--Weyl basis with entries $(-1)^{2j_e}$. It commutes with every Casimir
and with every original endpoint gauge action, because $g(-U)h^{-1}=-(gUh^{-1})$.
It fixes $\Omega_0$ and maps $W_f$ to $-W_f$ for every plaquette, selected or
omitted. Therefore
\begin{equation}
 U_E\,H_N(\tau,\kappa)\,U_E^*=H_N(-\tau,-\kappa),\qquad
 U_E\,Q_LH_N(\tau,\kappa)Q_L\,U_E^*=Q_L'\,H_N(-\tau,-\kappa)\,Q_L' ,
 \label{eq:aw1-flip-identity}
\end{equation}
in every open centered whole-star box. Here $Q_L'$ is the cutoff of
$h_b(-\kappa)$, which equals $Q_L$ at $\kappa=0$. For $\kappa\ne0$ the
identity rests on the inherited form of the strip operator: Casimirs, the
three selected face traces and scalars. A scalar shift cancels in the
ground-subtracted on-site operator. The reverse checker verifies
\eqref{eq:aw1-flip-identity} term by term on $\Lambda_2$: all on-site
Casimirs, the 1344 omitted and 375 selected face terms, at $\kappa=0$ and at
$\kappa=(\tfrac17,-\tfrac19,\tfrac13)$.

The contract also demands a positive demonstration at $\kappa\ne0$. The
forward producer compresses $H_0-(\tau/3)W-\kappa W_g$, with $g$ a selected
face, onto $\mathrm{span}\{\Omega_0,W\Omega_0,W_g\Omega_0\}$. The flip maps it
to $H(-\tau,-\kappa)$, not to $H(-\tau,\kappa)$. The reverse producer uses
two truncated plaquettes, one omitted face carrying $\tau$ and one selected
face carrying $\kappa=\tfrac17$. There the selected-face mean at
$(-\tau,-\kappa)$ is the exact negative of its value at $(\tau,\kappa)$, and
it is unchanged at $(-\tau,\kappa)$. Both are finite fixtures. Claiming
$\tau$-antisymmetry via $U_E$ at a nonzero triple is the rejected control
\recid{flip_breaks_at_nonzero_kappa}.

\paragraph{Covariance at the zero triple.}
At $\kappa=0$, AM2 gives a unique ground for both signs, in the full space and
in every cutoff space. Since $U_E\psi_N(\tau)$ is a ground of $H_N(-\tau)$,
$U_E\psi_N(\tau)\propto\psi_N(-\tau)$. With $U_E\Omega_0=\Omega_0$ and the
intermediate normalization $\langle\Omega_0,\psi\rangle=1$ this gives
$c_I(-\tau)=U_Ec_I(\tau)$, so $c^{(1)}$ is odd and $\norm{c}_a$ is even. Also
$\chi_{-\tau}=-U_E\chi_\tau$ and
$G_N(-\tau)-E_N=U_E\bigl(G_N(\tau)-E_N\bigr)U_E^*$. Hence, for every box and
every cutoff,
\begin{equation}
 \begin{gathered}
 \omega_{N,-\tau}=\omega_{N,\tau}\circ\alpha_E,\qquad
 \omega_{N,-\tau}(W)=-\omega_{N,\tau}(W),\\
 \omega_{N,-\tau}(W^2)=\omega_{N,\tau}(W^2),\qquad
 C_{N,-\tau}=C_{N,\tau},\qquad c_{N,-\tau}=c_{N,\tau}.
 \end{gathered}
 \label{eq:aw1-covariance}
\end{equation}

\paragraph{AQ limits as whole sets.}
Let $\mathcal S(\tau)$ be the set of all local trace-norm limits of
$\omega_{N_k,\tau}$ along subsequences of centered boxes. Conjugation by
$U_{E\cap F}$ on each finite region $F$ defines a consistent automorphism
$\alpha_E$ of the quasi-local algebra. Along any subsequence on which
$\omega_{N_k,\tau}$ converges, $\omega_{N_k,-\tau}$ converges to
$\omega_\tau\circ\alpha_E$. Therefore
\begin{equation}
 \mathcal S(-\tau)=\mathcal S(\tau)\circ\alpha_E\qquad\text{as sets of subsequential limits.}
 \label{eq:aw1-aq-sets}
\end{equation}
The automorphism $\alpha_E$ intertwines the finite-volume dynamics, and hence
their norm limits in AQ1. So for each $\omega\in\mathcal S(\tau)$,
$(\omega\circ\alpha_E)(W)=-\omega(W)$, while $\omega\circ\alpha_E$ has the same
$\omega(W^2)$, $c(\theta)$ and $C(s)$ as $\omega$. This is the AQ-level
evenness used in Section~\ref{sec:aw1-parity}. AQ1's chosen states at $+\tau$
and $-\tau$ come from $\tau$-dependent diagonal extractions. They are related
pointwise only along a common subsequence. In the reverse producer's toy
sequence the finite-volume identity holds exactly and the two limit sets are
negatives of each other, but independently chosen subsequences give the
nonzero sum $\tfrac1{500}$; the pointwise claim without a common subsequence
is rejected. Odd periodic sides and frozen or gauge-fixed links are excluded.

\paragraph{No third-order remainder from oddness.}
In each box, analyticity and oddness remove the $\tau^2$ Taylor coefficient
of $\omega_N(W)$, with a $\tau^3$ constant that depends on $N$. They do not
turn a uniform bound $|r|\le K_2\tau^2$ into $O(\tau^3)$: the function
$r(\tau)=K\tau|\tau|$ is odd and saturates the bound. The AQ limit, as a
function of $\tau$, is not even known to be continuous. Any third-order
remainder is a separate obligation, and every claim flag records none. A
loop-3 sign-off sentence that the direct remainder is $O(\tau^3)$ is used by
neither producer and is not admitted.

\paragraph{The modified-flip-set remark (recorded, not admitted).}
Both producers record, without claiming it, a second cochain. Every unit cube
contains $0$ or $2$ selected faces, so the selected-face indicator is a
$\mathbb Z_2$ cocycle and hence a coboundary $\delta E''$ on every box. The
forward producer's $E''$ consists of the $x$-links with $p_x\bmod4\in\{0,1,2\}$
and $p_y\bmod4\in\{1,2\}$; the reverse producer's is a period-8 cochain on $x$-
and $y$-links. The set $E^*=E\triangle E''$ then meets every omitted face
oddly and every selected face evenly, so $U_{E^*}$ would map
$H(\tau,\kappa)$ to $H(-\tau,\kappa)$ and make $\omega(W)$ odd in $\tau$ at
every selected triple. The skeptic confirmed the combinatorics independently:
both $E''$ have the right coboundary on a $17\times17\times7$ fine box, they
differ by a cocycle, and an exact $3\times3$ compression onto
$\{\Omega_0,2W\Omega_0,2W_g\Omega_0\}$ shows both identities. The operator
consequence at $\kappa\ne0$ is not reviewed. It needs four premises: the
on-site operator at a nonzero triple must contain only Casimirs, selected
face traces and scalars (the strip form inherited from A1, which neither
producer read and the skeptic did not review); the cutoffs must be spectral projections of $h_b(\kappa)$; the triple
must be held fixed independently of $\tau$; and a unique ground and AM2
applicability must hold with the non-Haar reference. The gate records the
remark and admits nothing at nonzero triples; a claim there would need its
own contract. The remark does not contradict the contract or the control
\recid{flip_breaks_at_nonzero_kappa}, both of which concern $U_E$. The
selection note's unqualified sentence ``antisymmetry in tau holds only at the
zero selected triple'' is to be read as ``via $U_E$''
(\repo{research/round32/skeptic/aw1.md}{review}, N3). The selection note is
historical and unchanged.

\subsection{The first-order Wilson mean}
\label{sec:aw1-mean}
\paragraph{Exact decomposition.}
Use the product split \eqref{eq:av1-order}--\eqref{eq:av1-orth}:
$\psi=\psi_{\rm out}+\delta$ with $\psi_{\rm out}=\Omega_R\otimes\varphi_{\rm out}$,
$n=\norm{\varphi_{\rm out}}$ and $e=\norm{\delta}/n$. Since $W$ acts on
$\mathcal H_R$ and $\langle\Omega_R,W\Omega_R\rangle=0$, and writing
$w=W\Omega_R$,
\begin{equation}
 \omega_N(W)=\frac{X+Y}{1+e^2},\qquad
 X=\frac{2\,\mathrm{Re}\langle w\otimes\varphi_{\rm out},\delta\rangle}{n^2},\qquad
 Y=\frac{\langle\delta,(W\otimes1)\delta\rangle}{n^2}.
 \label{eq:aw1-XY}
\end{equation}

\paragraph{Single-component overlap.}
The four links of $W$ have owners $0,0,e_z,0$, and $\norm{w}=\tfrac12$.
Integrating over the link $(0,x)$, owned by $0$, or over $(e_z,x)$, owned by
$e_z$, kills a single spin-$\tfrac12$ factor. So the partial vacuum
contraction $\langle\Omega_x|w\rangle$ vanishes for each $x\in R$, and $w$ is
excited at both sites of $R$. A creation meeting only one site of $R$ leaves
the vacuum at the other site and is orthogonal to $w$. This includes the 66
first-order straddling faces that meet $R$ once. At first order only the
creation supported exactly on $R$ overlaps $w$, with multiplier
$2\norm{w}=1$:
\begin{equation}
 -2\langle W\Omega_0,c^{(1)}_R\rangle
 =\frac{\tau}{36}\sum_{M_f=R}\mathbb E[WW_f]=\frac{\tau}{144}.
 \label{eq:aw1-overlap}
\end{equation}
Of the ten faces with owner set exactly $R$, only $f=W$ contributes. Hence
\begin{equation}
 \begin{gathered}
 \boxed{\;\omega_\tau(W)=-2\,\mathrm{Re}\langle W\Omega_0,c^{(1)}\rangle+O(\tau^2)
 =2\,\mathrm{Re}\langle W\Omega_0,\psi^{(1)}\rangle+O(\tau^2)
 =+\frac{\tau}{144}+r(\tau),\;}\\
 \omega^{(1)}_{\pm10^{-8}}=\pm\frac{1}{14400000000}.
 \end{gathered}
 \label{eq:aw1-mean}
\end{equation}
The expansion holds for either sign of $\tau$. Section~\ref{sec:aw1-k2}
bounds $r(\tau)$.

\paragraph{The sign chain.}
The sign is fixed by four steps, each from a frozen convention.
\begin{enumerate}[label=(\roman*)]
\item I1.5 puts $-(\alpha\tau/24)W_f$ into $H$ for every omitted face; in
normalized units the star is $\phi_b=-(\tau/3)\sum W_f$
\eqref{eq:dict-star}.
\item Each $W_f\Omega_0$ is an eigenvector of $H_0$ with energy $24$
(normalized) or $3$ (units of $\alpha$). The per-face coefficient of
$c^{(1)}=L_0=H_0^{-1}P_\perp V\Omega_0$ is therefore
$(-\tau/3)/24=(-\tau/24)/3=-\tau/72$ in both unit systems.
\item AM2 writes $\psi=e^{-C}\Omega_0=\prod_I(1-\hat c_I)\Omega_0$, so
$\psi^{(1)}=-c^{(1)}$, which is \eqref{eq:aw1-psi1}.
\item Only $f=W$ pairs with $W$, and $\mathbb E[W^2]=\tfrac14$, so the mean is
$2\cdot\tfrac{\tau}{72}\cdot\tfrac14=\tfrac{\tau}{144}$.
\end{enumerate}
For $\tau>0$ the interaction lowers the energy where $W>0$, and the mean is
positive. Mixing the unit systems gives $\tau/1152$ or $\tau/18$; both are
rejected controls. Three further re-derivations agree. The forward producer
minimizes over $\mathrm{span}\{\Omega_0,W\Omega_0\}$ in normalized units,
where $\langle W\Omega_0,V\Omega_0\rangle=-\tau/12$ and
$\langle W\Omega_0,HW\Omega_0\rangle=6$. The trial energy
$6\varepsilon^2-(\tau/6)\varepsilon$ is minimal at $\varepsilon=\tau/72$, and
the mean is then $2\varepsilon\,\mathbb E[W^2]=\tau/144$. The skeptic repeats
the minimization in units of $\alpha$. Both producers also use a
one-plaquette fixture, which is a finite graph with
\recid{transfers_to_aq} false. It is written in the character basis
$\chi_j(U_f)$, with $W\chi_j=\tfrac12(\chi_{j-1/2}+\chi_{j+1/2})$, free
energy $4j(j+1)$ in units of $\alpha$ and coupling $-(\tau/24)W$. Its exact
Rayleigh--Schr\"odinger series,
\begin{equation}
 \begin{gathered}
 \langle W\rangle=\frac{\tau}{144}+0\cdot\tau^2-\frac{5\,\tau^3}{11943936}+0\cdot\tau^4
 +\frac{289\,\tau^5}{6604518850560}+\cdots,\\
 E_2=-\frac{\tau^2}{6912},\qquad
 \langle W^2\rangle=\frac14+\frac{7\,\tau^2}{331776}+\cdots,
 \end{gathered}
 \label{eq:aw1-plaquette}
\end{equation}
is stable under truncation. At $j\le\tfrac12$ the exact ground values of
$\langle W\rangle$ at $\tau=\pm10^{-8}$ are negatives of each other and carry
the sign of $\tau$.

\paragraph{The contract's sign shorthand.}
Item~3 of the contract writes $\omega_\tau(W)=2\langle W\Omega_0,c^{(1)}\rangle$.
With AV1's creation coefficient $c^{(1)}=L_0$, the literal display evaluates
to $-\tau/144$. The correct identity is the one in \eqref{eq:aw1-mean}; the
display holds only if $c^{(1)}$ is read as the vector correction $\psi^{(1)}$.
Both producers derive $+\tau/144$ with the full sign chain and record the
shorthand as a convention finding. Both reject the literal reading as a
damaging mutation: the forward as
\recid{plus_two_L0_convention_gives_minus_tau_over_144}, the reverse as
\recid{am2_sign_literal} and \recid{contract_formula_with_am2_c1}. The skeptic
recorded the same defect before production. The gate records it as a
non-blocking wording defect. The frozen contract is not amended, and neither
the gate nor this section uses the display.

\paragraph{Face controls.}
The pairing $\mathbb E[WW_f]$ vanishes for the 81 other faces meeting $R$;
counting all ten faces inside $R$ would give $10\tau/144$. The Wilson face is
the $xz$ class with $r=0$, $q=0$ among the 21 omitted classes anchored at
factor~$0$ (Table~\ref{tab:dict-classes}). The other twenty are kept apart
throughout. Nine share the owner set $R$; they lie in $c^{(1)}_R$, but
their pairing with $W$ is zero by parity. Six have owner sets strictly
containing $R$, namely $\{0,e_x,e_z\}$ (two faces) and $\{0,e_y,e_z\}$ (four
faces); they reach $\omega(W)$ only through the straddling term. Five have
owner sets meeting $R$ only at $0$; they reach it only through the
two-creation and density terms. Merging the Wilson face with the others,
excluding it from the omitted set (coefficient $0$) and counting it twice are
rejected. For orientation invariance, every omitted class anchored at $0$,
used as the observable, has first-order mean $+\tau/144$, because only the
class face pairs with itself. Reversal and a cyclic shift of the base point
leave $\tfrac12\tr U_f$ unchanged. A selected face used as the observable
gives $0$ at the zero triple.

\subsection{The second-order remainder}
\label{sec:aw1-k2}
\paragraph{Outside excitations.}
One lemma beyond AV1 is needed. The reverse producer identifies it as the
missing premise, and the forward producer proves it separately. For a site
$u$ outside $R$, let $C_u$ collect the creations of $\varphi_{\rm out}$ that
contain $u$. They pairwise overlap, so $C_u^2=0$ and
$\varphi_{\rm out}=\varphi'-C_u\varphi'$ with
$\varphi'=(P_u\otimes1)\varphi_{\rm out}$. Hence
\begin{equation}
 \norm{(Q_u\otimes1)\varphi_{\rm out}}=\norm{C_u\varphi'}\le t_u\norm{\varphi'}
 \le t\,\norm{\varphi_{\rm out}},\qquad
 t_u=\sum_{I\ni u,\ I\cap R=\varnothing}\norm{c_I}\le t=\norm{c}_a .
 \label{eq:aw1-outside}
\end{equation}

\paragraph{Itemized bound.}
Let $a=|\tau|/144$ be the norm of one first-order face vector, and let $T$ and
$\varrho=352\,J\,T$ be the exact-tier quantities of
\eqref{eq:av1-selfconsistent}. Put $s_{\rm str}=6a+\varrho$, the sum over the
six first-order faces whose owner sets strictly contain $R$ ($6=16-10$), and
$s_0=s_{e_z}=33a+\varrho$, the sum over the faces whose owner sets meet $R$
at one site only ($33=49-16$). Inserting \eqref{eq:av1-delta} into
\eqref{eq:aw1-XY} and sorting the supports gives, with $\varepsilon\ge e$,
\begin{equation}
 \begin{aligned}
 \Bigl|\omega_N(W)-\frac{\tau}{144}\Bigr|\le{}&
 \underbrace{2\norm{w}\varrho}_{\recid{am2_remainder}}
 +\underbrace{2\norm{w}\,T\,s_{\rm str}}_{\recid{straddling}}
 +\underbrace{2\norm{w}\,s_0s_{e_z}}_{\recid{two_creation}}\\
 &+\underbrace{\norm{W}\varepsilon^2}_{\recid{density}}
 +\underbrace{\mathcal N}_{\recid{normalization_order}}
 =:K_2\tau^2 .
 \end{aligned}
 \label{eq:aw1-k2}
\end{equation}
The preregistered items are the following.
\begin{itemize}
\item \recid{am2_remainder}: $|2\mathrm{Re}\langle w,(c-c^{(1)})_R\rangle|
\le2\norm{w}\norm{(c-c^{(1)})_R}\le\varrho$, about $3354.108\,\tau^2$. It
dominates every exact-tier value.
\item \recid{overlap_multiplier}: $2\norm{W\Omega_R}=1$, with
$\norm{W\Omega_R}=\tfrac12$ a positive exact control. The factor $4$
($2\,\mathrm{Re}$, $\norm{W}=1$, two sites) is admitted only as a labelled
conservative variant, about $13416.96$ (forward) and $13416.82$ (reverse).
\item \recid{straddling}: only supports strictly containing $R$ pair with an
outside excitation, of amplitude at most $t\le T$ by \eqref{eq:aw1-outside}.
The 72 straddling faces split as $66+6$ by the size of $I\cap R$, and the 66
meeting $R$ once contribute exactly zero.
\item \recid{two_creation}: disjoint pairs $I\ni0$, $J\ni e_z$ contribute at
most $2\norm{w}A_0A_{e_z}$ with $A_x\le33a+\varrho$. This is conservative:
every first-order support has at least two sites, so the pair term is
actually of third order.
\item \recid{density}: $|Y|\le\norm{W}e^2\le\varepsilon^2$, from the traced
density $\tr_{\rm out}|\delta\rangle\langle\delta|$.
\item \recid{normalization_order}: the correction
$-(X+Y)e^2/(1+e^2)$ is of third order and is still charged. The leading term
is exactly $\tau/144$, so $a\varepsilon^2$ suffices; the reverse producer
charges the looser $(\varepsilon+\varepsilon^2)\varepsilon^2$.
\item \recid{arithmetic}: no numeric cost. Every value is an exact rational,
with triangle bounds instead of square roots.
\end{itemize}
The pins come from the contract enumeration: 15 owner sets with
multiplicities $\{1\times5,2\times3,3\times2,4\times3,10\times2\}$, 49 faces
per factor, 82 faces meeting $R$ and 16 touching both of its factors
(Table~\ref{tab:av1-owners}). Both producers derive them from the snapshotted
I1 table, and the skeptic re-derived $49$, $10$, $6$, $72=42+30$ and $33+33$.

\paragraph{Three valid values.}
The three itemizations differ in two places. For the density term, the
forward producer and the skeptic use AV1's $\varepsilon=2T+T^2$. The reverse
producer pins $\varepsilon$ to the 82 faces meeting $R$,
$\varepsilon_{\rm rev}=82a+2\varrho+s_0s_{e_z}$, which is AV1's labelled
82-face refinement applied inside the density term; it is sharper because
$2T\approx98a$ exceeds $82a$. For the straddling and pair terms, the skeptic's
pre-comparison headline charged $T\cdot T$ and $T^2$ without the counts.
Table~\ref{tab:aw1-k2} lists the terms. The skeptic also computed a minimal
accepted form that takes the pinned terms, the smaller $\varepsilon$ and
$a\varepsilon^2$. Every term of each itemization is at least the
corresponding minimal term, so all three are valid upper bounds, ordered
skeptic $>$ forward $>$ reverse. The skeptic's itemization coincides
exactly with the forward producer's labelled variant
\recid{unpinned_t_bounds}.

\begin{table}[htbp]
\centering\small
\caption{Exact-tier ledger of $K_2$ at $\tau=10^{-8}$, each entry divided by
$\tau^2$. Decimals are rounded previews of exact rationals in the frozen
outputs and the review record. A formula is shown where it differs from the
forward one.}
\label{tab:aw1-k2}
\begin{tabular}{llrrr}
\toprule
Term & Order & Forward & Reverse & Skeptic\\
\midrule
\recid{am2_remainder}, $\varrho$ & 2 & $3354.10836$ & $3354.10836$ & $3354.10836$\\
\recid{straddling} & 2 & $T(6a+\varrho)$: $0.014191$ & $0.014191$ & $T\cdot T$: $0.115812$\\
\recid{two_creation} & 2 & $(33a+\varrho)^2$: $0.052533$ & $0.052533$ & $T^2$: $0.115812$\\
\recid{density} & 2 & $(2T+T^2)^2$: $0.463247$ & $\varepsilon_{\rm rev}^2$: $0.324343$ & $0.463247$\\
\recid{normalization_order} & 3 & $a\varepsilon^2$: $3.2\times10^{-11}$ & $(\varepsilon+\varepsilon^2)\varepsilon^2$: $1.8\times10^{-9}$ & $3.2\times10^{-11}$\\
\recid{overlap_multiplier} & -- & 1 & 1 & 1\\
\recid{arithmetic} & -- & exact & exact & exact\\
\midrule
$K_2^+$ & & $3354.63832966$ & $3354.49942587$ & $3354.80322946$\\
\bottomrule
\end{tabular}
\end{table}

The gate binds the largest of the three as the admitted constant:
\begin{equation}
 \begin{aligned}
 K_2^+&=\frac{81108864767825329926713064490531229475390625}{24176936535511801466930759024724079017984}
 \approx3354.80323,\\
 K_2^+&\le\frac{3354803229461691}{10^{12}} .
 \end{aligned}
 \label{eq:aw1-K2}
\end{equation}
It uses only $T$, $\varrho$ and $\varepsilon=2T+T^2$, with no 6-, 33- or
82-face refinement. It therefore stays valid even if every refinement were
disputed. The forward value is exact; the reverse value is a rational with a
114-digit numerator, quoted here by its outward $10^{-40}$ ceiling:
\begin{equation}
 \begin{aligned}
 K_{2,\rm fwd}^+&=\frac{81104877995836618199905135109761286809699489}{24176936535511801466930759024724079017984}
 \approx3354.63833,\\
 K_{2,\rm rev}^+&\le\frac{33544994258669290845093439980837832702606211}{10^{40}}
 \approx3354.49943 .
 \end{aligned}
 \label{eq:aw1-K2-producers}
\end{equation}
The three values differ by at most about $0.304$.

The bound holds uniformly in $N$. Every constant depends only on $t$,
$\varrho$, $\varepsilon$ and bulk counts, finite-box counts never exceed
bulk counts, and the star of $W$ is retained for $N\ge2$. It holds at every
on-site cutoff $L\ge24$, where $c^{(1)}_L=c^{(1)}$. It passes to the
untruncated ground vector by AV1's cutoff removal
(Section~\ref{sec:av1-passage}) and, because $W$ is local on $R$, to every
AQ1 subsequential limit. Each $K_2(\tau)$ is nondecreasing in $|\tau|$, so
the cap value bounds every smaller coupling. The ratio $B(\tau)/B(\tau/100)$
of the total bounds is about $10000.98$, which confirms quadratic scaling. The
$-\tau$ ledger replays the same $|\tau|$ formula.

The AM2 generic majorant $352\,J\,T$ is $99.98\%$ of $K_2^+$. Parity would
remove the $\tau^2$ coefficient of the direct overlap
$\langle W\Omega_R,(c-c^{(1)})_R\rangle$, but using that needs a separate
third-order bound, which is neither used nor claimed. Sharper directed forms
of the AM2 remainder, of the type $G(t)-16\le288t/(1-8t)$, give labelled
variants of about $2745$ (forward $2744.88$, reverse $2744.66$). They are
recorded, not admitted.

\paragraph{The crude tier fails and is retained.}
The crude tier uses only $t_{\rm c}=J_0G(R_a)=592|\tau|=37/6250000$ at the
cap, with no face counts, so no tiers are mixed. Its terms divided by
$\tau^2$ are $5834752$ (AM2 remainder $352\,J\,t_{\rm c}$), $350464$
(straddling), $350464$ (two-creation) and about $1401864.30$ (density).
Together they give
\begin{equation}
 K_2^{\rm crude}=\frac{1937877026146766159414097129}{244140625000000000000}
 \approx7.9375443\times10^{6}
 \label{eq:aw1-K2crude}
\end{equation}
for the forward producer and the skeptic, and about $7.9375609\times10^6$ for
the reverse producer, which differs only in the normalization charge. Then
$K_2^{\rm crude}\tau\approx0.0794$, which exceeds $1/288$ by a factor of about
$22.9$. The sign margin is $0.0875$. The crude tier fails at the cap and is
retained as a limited-tier value. Its rule value $10^{-10}$ is information
only and is never substituted.

\subsection{The frozen AW2 coupling rule}
\label{sec:aw1-rule}
The contract freezes the rule before any constant is known: $\tau_{\rm AW2}$
is ``the largest element of the decade grid $\{10^{-8},10^{-9},\ldots\}$ with
$K_2^+\tau\le1/288$'', where $K_2^+$ is the admitted exact-tier constant.
Given $|r(\tau)|\le K_2^+\tau^2$ and $\tau>0$, the interval
$\tau/144\pm K_2^+\tau^2$ excludes zero with margin
$(\tau/144)/(K_2^+\tau^2)=1/(144K_2^+\tau)$. The threshold $1/288$ is exactly
a margin of $2$, the floor fixed by a loop-2 veto (Section~\ref{sec:plan}).
With the bound value \eqref{eq:aw1-K2}, $K_2^+\tau\approx3.35480\times10^{-5}$
at $\tau=10^{-8}$, which is at most $1/288$. The rule therefore gives
\begin{equation}
 \begin{gathered}
 \tau_{\rm AW2}=10^{-8}\ (\text{the cap}),\\
 \frac{1}{144K_2^+\tau}
 =\frac{42981220507576535941210238266176140476416}{207638693805632844612385445095759947457}
 \approx207.00005\ge\frac{41400010489}{200000000}.
 \end{gathered}
 \label{eq:aw1-rule}
\end{equation}
The choice among the three valid values changes no decision. Any exact-tier
constant up to $10^8/288\approx347222$ gives $10^{-8}$, and the forward and
reverse values give margins of about $207.010$ and $207.019$. The quotient of
$1/288$ by $K_2^+\tau$ is about $103.5$. The reverse headline ``margin
$\approx103.5$'' is that quotient, not the contract's sign margin.

The admitted numbers are compatible with AV1. The tier-(ii) bound
\eqref{eq:av1-Dii} satisfies
$D_{\rm ii}\approx1.3612\times10^{-8}\ge\tau/144+K_2^+\tau^2\approx6.978\times10^{-11}$,
about 196 times the first-order term. The feasibility Boolean is not a sign
certificate. AW1 admits no enclosure and no sign of $\omega(W)$; instantiating
the interval at $\tau_{\rm AW2}$ is AW2's task. The gate advises that AW2's
checker bind $K_2^+$ to the rational \eqref{eq:aw1-K2} from the gate rather
than recompute it.

\subsection{Controls, replays and review}
\label{sec:aw1-evidence}
The contract lists 30 controls. The forward checker runs 41 exact checks and
rejects 68 damaging mutations; the reverse checker runs 57 exact checks and
rejects 107. Each implements all 30 controls as damaging mutations, and
neither uses a floating-point value or \texttt{sympy} in the admission path.
Both record the same claim flags:
\begin{itemize}\raggedright
\item true: \recid{first_order_parity_claim} and \recid{flip_lemma_claim};
\item false: \recid{third_order_remainder_claim}, \recid{continuum_claim},
\recid{uniform_wilson_claim}, \recid{resolved_interaction_shift} and
\recid{scientific_priority_verified}.
\end{itemize}
Both checkers verify the
snapshot's SHA-256 before parsing and read the target $1/288$, the rule and
the reference values $0$ and $\tfrac14$ from it. After a coherently rehashed
edit of the target to $1/144$, the reverse checker rejects the target and the
forward checker records it. That edit requires changing \texttt{check.py},
which the freeze inventory binds, so the gate records the difference as
non-blocking.

The skeptic's pre-comparison package, with 79 exact checks, was committed
before either producer's commit. The post-comparison review adds 65 exact
checks, including 30 source-edit mutation receipts. Twenty-five of the edits
abort at the intended check, and one shows that a value is read from the
contract. Four silent value edits run to completion in the producer
checkers: forward density divided by 4, forward pair term divided by
9, reverse density pinned to 49 faces, and reverse pair term divided by 9.
Only the skeptic's term-by-term validator catches them, because neither
producer pins individual term values. All four producer replays, forward and
reverse under normal and optimized Python, reproduce the frozen outputs byte
for byte. The closures hold 39 and 35 files, and the reverse inputs are
exactly \texttt{AGENTS.md}, the contract and the 29 shared premises.

One isolation finding is recorded without being resolved. The shared session
scratchpad root holds files written between 23:14Z and 23:20Z, before either
producer commit. Their content matches the forward producer's $K_2$ formulas,
its $E''$ remark and its one-plaquette fixture, and the reverse producer's
$\alpha$-unit fixture and four-site creation fixture. Neither producer
disclosed scratch reads. The forward producer's file overwrote the
skeptic's pre-comparison scratch file of the same name. Whether any agent read
another's scratch file cannot be determined. The routes and code differ
materially: two different $E''$ constructions, two different density
$\varepsilon$ and different fixtures. Independence is therefore not refuted,
but premise isolation is verified for repository inputs only.

\subsection{Limitations recorded by the gate}
\label{sec:aw1-limits}
The AW1 gate records the following limitations verbatim.
\begin{gatequote}
\begin{itemize}
\item Only the AM2/AQ1 zero-selected patterned family, the cover R=\{0,e\_z\},
fixed spacing and |tau|<=10\textasciicircum{}-8; nothing transfers to nonzero
selected triples, uniform Wilson theory, weak coupling or the continuum
\item Parity theorem for C(s) and c(theta) is a finite-volume Taylor statement
(Kato analyticity per box); the tau\textasciicircum{}2 constant is explicitly
unbounded; AQ limits inherit only the set-level evenness
\item Flip lemma needs open boxes and kappa=0 for antisymmetry; the kappa!=0
identity rests on the inherited strip-operator form; AQ statements are about
whole sets of subsequential limits
\item K\_2\textasciicircum{}+ is an absolute second-order upper bound with no
sign of r(tau) and no third-order remainder; the AM2 generic majorant 352JT is
99.98\% of it; the -tau ledger is a replay
\item No enclosure or sign of omega(W) is admitted (AW2); static equal-time
mean, not a dynamical or mass-gap correction
\item Inherited without re-proof: AM2 fixed point, majorant, gaps and ground
uniqueness; AV1 product split, anchored-norm tier, cutoff removal and AQ
passage; AQ1 construction and dynamics; I1 dictionary; Kato and Peter-Weyl
cited, not machine-checked
\item Fixtures (one plaquette, qubit/four-site creation algebras, kappa
compressions, 3x3 compression) are finite graphs or algebras,
transfers\_to\_aq false
\item Independence limited to routes and code: shared premises state the
mechanism (loop-3 sign-offs, selection-aw1.md, skeptic/av2.md); the forward
read the skeptic triage and loop-2 notes as declared premises
\item N1: both producers used the shared scratchpad root (files matching the
forward\textquotesingle{}s K\_2/K\_W2 formulas,
E\textquotesingle{}\textquotesingle{} remark and one-plaquette fixture, and
the reverse\textquotesingle{}s alpha-unit fixture and four-site creation
fixture, 23:14-23:20Z) without disclosing scratch reads; the
forward\textquotesingle{}s k2.py overwrote the skeptic\textquotesingle{}s
pre-comparison scratch k2.py; reads cannot be determined; isolation is
verified for repository inputs only
\item N4: the reverse\textquotesingle{}s headline \textquotesingle{}margin
\textasciitilde{}103.5\textquotesingle{} is the ratio to 1/288; the contract
sign margin is 1/(144 K\_2\textasciicircum{}+ tau) \textasciitilde{}207.02
\item N5: forward K\_W2 \textasciitilde{}3354.57 is conservative (2 rho, all
72 straddling faces) against reverse 1677.33 and skeptic 1677.52;
supplementary, not a target
\item N6: reverse wording that U\_E commutes with every on-site spectral
cutoff and fixes P\_R holds at kappa=0; at kappa!=0 these map to the -kappa
objects (forward states Q\_L\textquotesingle{})
\item N7: the forward records a coherently rehashed 1/144 target instead of
rejecting it (no semantic tie to half the coefficient; the reverse rejects);
protected by the check.py hash in freeze.json
\item N8: neither producer checker pins individual K\_2 term values; four
silent undercounts ran to completion and were caught only by the skeptic term
validator; AW2 should bind K\_2\textasciicircum{}+ from the gate
\item N9: the modern loop-3 sign-off sentence that the direct remainder is
O(tau\textasciicircum{}3) is used by neither producer and is not admitted
\item Undeclared disclosed reads: forward tools README, freeze.py, AV1 forward
check.py and two assistant scripts (named, not imported; the modern K\_2
script holds only loop-2 tiers); reverse portions of reverse/av1/check.py and
the line count of reverse/av2/check.py; the reverse\textquotesingle{}s partial
reads of the AV2 forward report and skeptic/av2.md are of declared shared
premises
\item Scientific priority unverified
\end{itemize}
\end{gatequote}
Full derivations and reviewed implementations:
\repo{research/round32/forward/aw1/report.md}{AW1 forward},
\repo{research/round32/reverse/aw1/report.md}{AW1 reverse},
\repo{research/round32/skeptic/aw1.md}{independent model-agent review} and the
\repo{research/round32/contracts/aw1.json}{frozen contract}. The gate selects
AW2 next.

\subsection{AW2: sign-certified enclosure}
\label{sec:aw2}
To be completed from the AW2 gate.
